\section{Performance Analysis}
\label{sec:performance}

\subsection{GerryChain Baseline Measurements}

We measured \gc{} throughput by running 1{,}000 \recom{} steps for
three states and recording wall-clock time.
These are direct empirical measurements from the \gc{} Python package
(version 0.3.2, NumPy 1.24, Python 3.10) run on an Intel Core i7-11800H
(2.3~GHz base, 4.6~GHz boost, 24~MB L3 cache) under Windows~11.
The starting plan for each state is a cold start from a seed-0 random
valid plan.
No warm-up steps were excluded; all 1{,}000 steps count toward the
wall-clock total.
\gc{} was run without pair reselection (the standard configuration):

\begin{center}
\begin{tabular}{lrrrl}
\toprule
State & $k$ & $n_{\text{tracts}}$ & 1K steps (sec) & Steps/sec \\
\midrule
North Carolina & 14 & 2{,}672 & 47 & $\approx 21$ \\
Wisconsin      &  8 & 1{,}922 & 28 & $\approx 36$ \\
Georgia        & 14 & 1{,}978 & 30 & $\approx 33$ \\
\bottomrule
\end{tabular}
\end{center}

The variation in throughput across states reflects differences in both
$n$ (tract count) and $k$ (district count).
Wisconsin is faster than North Carolina despite having fewer districts
($k=8$ vs $k=14$) because its lower tract count reduces the merged-region
size; Georgia achieves similar throughput to Wisconsin despite $k=14$
because its tracts are on average smaller in the adjacency graph.

\subsection{Rust Target Estimates}

The theoretical throughput of \re{} follows from Wilson's complexity
analysis.
For a merged region of size $m = 2n/k$:

\begin{enumerate}
  \item \textbf{Wilson's walk}: $O(m\log m)$ expected operations.
    At $m \approx 191$ (NC average: $2 \times 2{,}672 / 14$),
    Wilson requires approximately $191 \times \log(191) \approx 191 \times 7.5 \approx 1{,}433$
    random walk steps.
  \item \textbf{Effective operations per walk step}: Neighbor selection
    (1 uniform draw), loop detection (hash set lookup), and
    path extension (pointer write)---approximately 8 instructions at the
    clock level.
  \item \textbf{Total operations per \recom{} step}: $\approx 191 \times 8 = 1{,}528$
    effective operations.
  \item \textbf{Effective clock rate}: Modern CPUs execute at
    $\sim$4\,GHz but MCMC-adjacent code (random number generation,
    irregular memory access) achieves roughly 1\,GHz effective throughput
    on graph-walk workloads of this size.
  \item \textbf{Resulting time per step}: $1{,}528 \,/\, 10^9 \approx 1.5\,\mu$s.
  \item \textbf{Steps per second}: $\approx 660{,}000$.
\end{enumerate}

Applying a conservative $13\times$ overhead factor for Rust runtime costs
(subgraph construction, adjacency list traversal, SmallVec allocation,
serde bookkeeping) yields a practical target of \textbf{50{,}000 steps/sec}.
This is still a $2{,}350\times$ speedup over \gc{} for North Carolina.

\begin{remark}
The overhead factor of $13\times$ is deliberately conservative.
Our METIS pipeline achieved a $213\times$ speedup over Python for a
structurally similar graph-partitioning workload~\citep{dellaLibera2026b1}.
The overhead factor for \re{} is higher because Wilson's random walk has
more irregular memory access patterns than METIS's multilevel contraction,
but $50{,}000$\,steps/sec should be achievable without architecture-specific
optimization.
\end{remark}

\paragraph{Overhead sensitivity.}
The $13\times$ overhead factor is an estimate; Table~\ref{tab:overhead}
shows how the projected throughput and speedup change under three scenarios.

\begin{table}[h]
\centering
\caption{Throughput sensitivity to overhead factor (NC baseline: 21 steps/sec).}
\label{tab:overhead}
\begin{tabular}{lrrr}
\toprule
Scenario & Overhead factor & Steps/sec & Speedup vs.\ \gc{} (NC) \\
\midrule
Optimistic  & $6.5\times$  & $\sim$100{,}000 & $\sim$4{,}700$\times$ \\
Central     & $13\times$   & $\sim$50{,}000  & $\sim$2{,}350$\times$ \\
Conservative & $26\times$  & $\sim$25{,}000  & $\sim$1{,}200$\times$ \\
\bottomrule
\end{tabular}
\par\vspace{0.5ex}
\small All figures estimated; Phase~2 \texttt{criterion.rs} benchmarks will determine the
actual overhead factor and narrow the range.
\end{table}

\subsection{Projected Performance Table}

\begin{table}[h]
\centering
\caption{%
  Performance comparison: \gc{} measured vs.\ \re{} estimated.
  Rust estimates based on Wilson's $O((n/k)\log(n/k))$ analysis at
  50{,}000\,steps/sec (conservative).
  PA and TX estimated from $n$-scaling of the NC/WI/GA measurements.
  All Rust figures require validation against \texttt{criterion.rs}
  (Phase~2).
}
\label{tab:performance}
\begin{tabular}{lrrrrrr}
\toprule
\multirow{2}{*}{State} & \multirow{2}{*}{$k$} &
  \multirow{2}{*}{$n$} &
  \multicolumn{2}{c}{\gc{} (measured)} &
  \multicolumn{2}{c}{\re{} (estimated\est)} \\
\cmidrule(lr){4-5}\cmidrule(lr){6-7}
& & & 1K steps (s) & Steps/s & 10K steps (s) & Speedup \\
\midrule
NC & 14 & 2{,}672 & 47 & 21 & 0.20 & 2{,}350$\times$ \\
WI &  8 & 1{,}922 & 28 & 36 & 0.12 & 2{,}330$\times$ \\
GA & 14 & 1{,}978 & 30 & 33 & 0.13 & 2{,}310$\times$ \\
PA & 17 & 3{,}218 & $\sim$30\est & $\sim$33 & 0.18 & 2{,}200$\times$ \\
TX & 38 & 4{,}866 & failure\bench & --- & 0.31 & ---\bench \\
CA & 52 & 8{,}057 & failure\bench & --- & 0.52 & ---\bench \\
\bottomrule
\end{tabular}
\par\vspace{0.5ex}
\small
\est\,Rust estimates require validation with \texttt{criterion.rs}; see Section~\ref{ssec:phase2}.\\
\bench\,\gc{} \emph{without} pair reselection fails on TX ($k=38$) and CA ($k=52$) from cold start
due to bipartition infeasibility; \gc{} \emph{with} pair reselection and \re{} both handle these cases
via pair reselection (Section~\ref{ssec:tx}).
\gc{}'s bipartition failures on TX $k=38$ are combinatorial, not language-specific.
Both GerryChain-with-pair-reselection and \re{} succeed.
The advantage of \re{} is throughput, not correctness.
\end{table}

\subsection{Texas and California}
\label{ssec:tx}

The bipartition infeasibility problem in Texas is a combinatorial property
of the TX plan space, not a language-specific limitation.
\gc{} without pair reselection is documented to fail on Texas from cold
start~\citep{gerrychain2021}; \gc{} with pair reselection would also
handle these states, albeit more slowly.
\re{}'s advantage is that pair reselection at microsecond speed (rather
than second speed in Python) reduces warm-up wall time without altering
the algorithm's combinatorial behavior.
The performance estimate in Table~\ref{tab:performance} assumes the pair
reselection overhead is absorbed within the 50{,}000\,steps/sec budget.

Per-step cost is $O((n/k)\log(n/k))$, so the relevant quantity is the
average merged-region size $m \approx 2n/k$.
For CA ($k=52$, $n=8{,}057$): $m \approx 310$.
For PA ($k=17$, $n=3{,}218$): $m \approx 379$.
PA has the largest average merged region among the six G-track states
and therefore the highest per-step cost; California's per-step cost
is lower than PA's despite its larger tract count, because its high $k$
reduces the merged-region size.
\emph{The actual per-step cost ordering across all six states is deferred
to Phase~2 empirical benchmarks.}
The estimated 0.52 seconds for 10{,}000 CA steps and 0.18 seconds for
10{,}000 PA steps (Table~\ref{tab:performance}) reflect the theoretical
ordering; Phase~2 criterion.rs results will confirm or correct the
estimates.

\subsection{Phase 2: Criterion.rs Benchmarks}
\label{ssec:phase2}

\emph{All Rust throughput figures in this paper are estimates derived from
complexity analysis and are deferred to Phase~2 empirical validation.}
Phase~2 will validate these estimates using \texttt{criterion.rs}, Rust's
standard micro-benchmarking harness, with the following protocol:

\paragraph{Hardware.}
Benchmarks will be run on the same Intel Core i7-11800H workstation used
for the \gc{} baseline measurements in Section~\ref{sec:performance}
(2.3~GHz base, 4.6~GHz boost, 24~MB L3 cache, Windows~11, single-threaded
to isolate per-step cost).
The CPU model, clock speed, and OS will be reported with results.

\paragraph{GerryChain comparison.}
The Phase~2 comparison will use \gc{} version~0.3.2, NumPy~1.24,
Python~3.10, starting from a cold-start seed-0 random valid plan for
each state (matching the Section~\ref{sec:performance} baseline).
Wall-clock time will be measured for 1{,}000 steps after a 50-step
warm-up (the warm-up steps are excluded from the throughput count).

\paragraph{Acceptance criterion.}
The Phase~2 report will confirm the 50{,}000\,steps/sec estimate if the
\texttt{criterion.rs} measurement for NC exceeds 30{,}000\,steps/sec
(i.e., within 40\% of target).
If the measurement falls below 30{,}000\,steps/sec, the abstract and
Table~\ref{tab:performance} speedup figures will be revised downward.
Results will be reported as 95\% confidence intervals from
\texttt{criterion.rs}, not point estimates.

\paragraph{Granularities.}
Benchmarks will be run at the following levels:
\begin{itemize}
  \item \textbf{Wilson UST}: Time to generate one uniform random spanning
    tree for a merged region of size $m$, as a function of $m$.
  \item \textbf{Balance check}: Time for the full edge enumeration pass.
  \item \textbf{Full step}: End-to-end \recom{} step time including
    subgraph construction.
  \item \textbf{Chain throughput}: Steps/sec for a 1{,}000-step chain,
    compared directly against \gc{} wall-clock time on the same hardware.
\end{itemize}

Results will be reported in a companion note to this paper.
